\section{Configurazione delle cartelle e del sistema}
\subsection{Struttura generale}
Il sito web è organizzato in cartelle e sottocartelle per mantenere una struttura ordinata e facilitare la gestione dei file.\\
\textbf{Le principali cartelle sono:}
\lstset{
    language=Python,
    basicstyle=\ttfamily\fontsize{11pt}{13pt}\selectfont,
    backgroundcolor=\color{gray!10},
    frame=single,
    breaklines=true,
    keywordstyle=\color{blue}\bfseries,  % grassetto + blu
    morekeywords={src,assets,database,relazione,index,htaccess},
    commentstyle=\color{gray!80},
    showstringspaces=false
}


\begin{lstlisting}[language=Python]
assets/
|-- icons/         # loghi, icone
|-- images/...     # foto di animali, di utenti, di admin ecc.

database/
|-- schema.sql     # schema del database
|-- populate.sql   # script per il popolamento iniziale del database (già eseguito sul server)

src/
|-- css/           # file per stili con pagine esterne: suddivisione per utente e admin, breakpoints e stampa 
|-- js/            # file JavaScript minimo per abbellimenti
|-- php/           # file PHP per la logica di creazione delle pagine e logica interna alle stesse.
|-- html/
    |-- main/      # file statici contenuto principale statico (quello racchiuso solitamente in <main>...</main>).
    |-- layout.html        # layout per utenti
    |-- layout-admin.html  # layout per admin
|-- DBconnection.php  # file per la connessione al database e le funzioni per eseguire query
|-- utils.php         # file con funzioni di utilità generali

relazione/
|-- ...            # file della relazione (questo file)

index.php          # file di routing principale del sito web
.htaccess         # file di configurazione del server web Apache
\end{lstlisting}
Le cartelle main e php raccolgono le pagine di admin in una sottocartella dedicata (\texttt{admin}), per chiarezza e organizzazione.

\subsection{file .htaccess e index.php di routing}
\subsection{Come avviene la costruzione delle pagine}